\documentclass{article}
\usepackage[margin=1in]{geometry}
\usepackage{amsmath, amssymb, amsthm, enumitem, multicol}

\everymath{\displaystyle}
\newtheorem*{theorem*}{Theorem}
\newcommand{\vct}[1]{\langle #1 \rangle}
\newcommand{\vu}{\mathbf{u}}
\newcommand{\vv}{\mathbf{v}}
\newcommand{\vw}{\mathbf{w}}
\newcommand{\vx}{\mathbf{x}}
\newcommand{\vn}{\mathbf{n}}

\newcommand{\colvecTwo}[2]{
  \begin{bmatrix}
    #1 \\
    #2
  \end{bmatrix}
}
\newcommand{\colvecThree}[3]{
  \begin{bmatrix}
    #1 \\
    #2 \\
    #3
  \end{bmatrix}
}

\newcommand{\casesTwo}[2]{
  \begin{cases}
    #1 \\
    #2
  \end{cases}
}
\newcommand{\casesThree}[3]{
  \begin{cases}
    #1 \\
    #2 \\
    #3
  \end{cases}
}

\title{Linear Algebra: A Modern Introduction}
\author{David Poole}
\date{}

\begin{document}
\maketitle
\section{Vectors}
\subsection{The Geometry and Algebra of Vectors}
\begin{enumerate}
    \setcounter{enumi}{3}
  \item
    \begin{enumerate}
      \item (3, 0, 1)
      \item (0, 0, 0)
      \item (2, 4, 0)
      \item (4, 3, 3)
    \end{enumerate}
    \setcounter{enumi}{12}
  \item $\mathbf{u}=(\frac{1}{2}, \frac{\sqrt{3}}{2})$,
    $\mathbf{v}=(-\frac{\sqrt{3}}{2}, -\frac{1}{2})$,
    $\mathbf{u}+\mathbf{v}=(\frac{1-\sqrt{3}}{2}, \frac{\sqrt{3}-1}{2})$,
    $\mathbf{u}-\mathbf{v}=(\frac{1+\sqrt{3}}{2}, \frac{\sqrt{3}+1}{2})$
  \item
    \begin{enumerate}
      \item $\mathbf{b}-\mathbf{a}$
      \item $-\mathbf{a}$
      \item $-2\mathbf{a}$
      \item $2(\mathbf{a}-\mathbf{b})$
      \item $\mathbf{b}-2\mathbf{a}$
      \item $\mathbf{a}-\mathbf{a}$
    \end{enumerate}
    \setcounter{enumi}{20}
  \item $-2\mathbf{u}+4\mathbf{v}$
    \setcounter{enumi}{39}
  \item 1
    \setcounter{enumi}{43}
  \item 4
  \item 2
  \item 2
  \item No solution.
  \item 3
  \item 3
  \item No solution.
  \item No solution.
  \item 8
  \item 2
    \setcounter{enumi}{56}
  \item
    \begin{enumerate}
      \item $a=\{1, 2, 3, 4\}$
      \item $a=\{1, 5\}$
      \item $a$ and $m$ can have no common factors other than 1 (i.e., the greatest common
        divisor ($\gcd$) of $a$ and $m$ is 1).
    \end{enumerate}
\end{enumerate}

\subsection{Length and Angle: The Dot Product}
\begin{theorem*}[The Cauchy-Schwarz Inequality]
  For all vectors $\mathbf{u}$ and $\mathbf{v}$ in $\mathbb{R}^n$,
  \[|\mathbf{u}\cdot\mathbf{v}|\le\left\|\mathbf{u}\right\|\left\|\mathbf{v}\right\|.\]
\end{theorem*}
\begin{theorem*}[The Triangle Inequality]
  For all vectors $\mathbf{u}$ and $\mathbf{v}$ in $\mathbb{R}^n$,
  \[\lVert\mathbf{u}+\mathbf{v}\rVert\le\left\|\mathbf{u}\right\|+\left\|\mathbf{v}\right\|.\]
\end{theorem*}
\begin{proof}
  \begin{align*}
    \lVert\mathbf{u}+\mathbf{v}\rVert^2 & = (\mathbf{u}+\mathbf{v})\cdot(\mathbf{u}+\mathbf{v}) \\
    & = \mathbf{u}\cdot\mathbf{u} +
    2(\mathbf{u}\cdot\mathbf{v}) + \mathbf{v}\cdot\mathbf{v}                                    \\
    & \le \mathbf{u}\cdot\mathbf{u} +
    2\left|\mathbf{u}\cdot\mathbf{v}\right| + \mathbf{v}\cdot\mathbf{v}                         \\
    & \le \lVert\mathbf{u}\rVert^2 +
    2\lVert\mathbf{u}\rVert\lVert\mathbf{v}\rVert + \lVert\mathbf{v}\rVert^2                    \\
    & =
    \left(\lVert\mathbf{u}\rVert+\lVert\mathbf{v}\rVert\right)^2.
  \end{align*}
  Since both sides are non-negative, taking the square root gives,
  \[\lVert\mathbf{u}+\mathbf{v}\rVert\le\left\|\mathbf{u}\right\|+\left\|\mathbf{v}\right\|.\]
\end{proof}
\subsubsection*{Exercises}
\begin{enumerate}
    \begin{multicols}{6}
    \item -1
    \item 0
    \item 11
    \item 2.62
    \item 2
    \item 3.6265
    \end{multicols}
  \item $\sqrt{5}$; $\left[\frac{-1}{\sqrt{5}}, \frac{2}{\sqrt{5}}\right]$
  \item $\sqrt{13}; \left[\frac{3}{\sqrt{13}}, \frac{-2}{\sqrt{13}}\right]$
  \item $\sqrt{14}; \left[\frac{1}{\sqrt{14}}, \frac{2}{\sqrt{14}}, \frac{3}{\sqrt{14}}\right]$
    \setcounter{enumi}{12}
  \item $\sqrt{17}$
  \item $\sqrt{65}$
  \item $\sqrt{6}$
  \item $\approx5.14$
  \item
    \begin{enumerate}
      \item Taking the length (or the norm) of $\mathbf{u}\cdot\mathbf{v}$, which is a scalar
        doesn't make sense.
      \item $\mathbf{u}\cdot\mathbf{v}$ is a scalar which cannot be added to a vector.
      \item $\mathbf{v}\cdot\mathbf{w}$ is a scalar, which does not have a dot product operation
        with a vector.
      \item $c$ is a scalar which does not have a dot product operation with $\mathbf{u}+
        \mathbf{v}$, which is a vector.
    \end{enumerate}
  \item Obtuse.
  \item Acute.
  \item Right angle.
  \item Acute.
  \item Obtuse.
  \item Acute.
    \setcounter{enumi}{24}
  \item $\frac{\pi}{3}$
    \setcounter{enumi}{26}
  \item $\approx1.5376$
    \setcounter{enumi}{28}
  \item $\approx0.2502$
    \setcounter{enumi}{30}
  \item $\overrightarrow{AB}$ is orthogonal to $\overrightarrow{AC}$.
  \item Angle between the diagonal $[1, 1, 1]$ and edge $[1, 0, 0]$ is $\approx0.9553$.
  \item Their dot products are not 0.
  \item The diagonals are perpendicular. The length of one side is $\frac{\sqrt{19}}{2}$.
  \item $(-2, 1, 1)$.
  \item $(200, -40)$
  \item $(4, 3)$
  \item
    \begin{enumerate}
      \item $0.5$ km
      \item 6 minutes
    \end{enumerate}
  \item Suppose Bert's velocity is $\mathbf{s}=\langle s_x, s_y \rangle$. We have
    $\|s\|=2$. To swim straight across the river with a flowing rate of $1$ mph, $s_x$
    needs to be $-1$. Solving for $s_y$, $\sqrt{s_x^2+s_y^2}=2 \implies s_y=\sqrt{3}$.
    Thus, the angle between $\mathbf{s}$ and the river bank is $60^\circ$.
  \item $\vct{-3, 3}$
  \item $\vct{\frac{-3}{5}, \frac{4}{5}}$
    \setcounter{enumi}{42}
  \item $\vct{3/2, -3/2, 3/2, -3/2}$
    \setcounter{enumi}{46}
  \item $\frac{3\sqrt{5}}{2}$
  \item $\frac{1}{5}$
  \item $\{3, 2\}$
  \item All vectors in standard position whose heads lie on the line $y=-3x$.
  \item All vectors in standard position whose heads lie on the line $y=-\frac{a}{b}x$.
  \item
    \begin{enumerate}
      \item $\theta=0$.
      \item $\theta=\pi$.
    \end{enumerate}
    \setcounter{enumi}{58}
  \item Applying the triangle inequality to vector $\vu-\vv+\vv$ gives
    \begin{align*}
      \|\vu-\vv+\vv\|           & \le \|\vu-\vv\|+\|\vv\| \\
      \|\vu-\vv+\vv\| - \|\vv\| & \le \|\vu-\vv\|         \\
      \|\vu\| - \|\vv\|         & \le \|\vu-\vv\|         \\
      \|\vu-\vv\|               & \ge \|\vu\| - \|\vv\|
    \end{align*}
  \item No. For example, $\vu = \langle1, 1\rangle, \vv = \langle1, -1\rangle,
    \vw = \langle2, -2\rangle$.
    \setcounter{enumi}{65}
  \item 9
  \item Given that $\vu\cdot\vv=3$, we can assume that they are non-zero vectors. Consider the
    vector $\vu-\vv$, it should be the case that $\|\vu-\vv\|>0$. However, $\|\vu-\vv\|=
    \|\vu\|^2-2\vu\cdot\vv+\|\vv\|^2=1-6+4=-1<0$, a contradiction. Thus, such vectors do not
    exist.
\end{enumerate}

\subsection{Lines and Planes}
\begin{enumerate}
  \item $\colvecTwo{3}{2}\cdot\colvecTwo{x}{y}=\colvecTwo{3}{2}\cdot\colvecTwo{0}{0}$, $3x+2y=0$
  \item $\colvecTwo{3}{-4}\cdot\colvecTwo{x}{y}=\colvecTwo{3}{-4}\cdot\colvecTwo{1}{2}$, $3x-4y=-5$
  \item $\vx=\colvecTwo{1}{0}+t\colvecTwo{-1}{3}$, $\casesTwo{x=1-t}{y=3t}$
  \item $\vx=\colvecTwo{-4}{4}+t\colvecTwo{1}{1}$, $\casesTwo{x=-4+t}{y=4+t}$
  \item
    \begin{enumerate}
      \item $\vx=\colvecThree{0}{0}{0}+t\colvecThree{1}{-1}{4}$
      \item $\casesThree{x=0+t}{y=0-t}{z=0+4t}$
    \end{enumerate}
  \item
    \begin{enumerate}
      \item $\vx=\colvecThree{3}{0}{-2}+t\colvecThree{2}{5}{0}$
      \item $\casesThree{x=3+2t}{y=5t}{z=-2}$
    \end{enumerate}
  \item
    \begin{enumerate}
      \item $\colvecThree{3}{2}{1}\cdot\colvecThree{x}{y}{z}=2$
      \item $3x+2y+z=2$
    \end{enumerate}
  \item
    \begin{enumerate}
      \item $\colvecThree{2}{5}{0}\cdot\colvecThree{x}{y}{z}=6$
      \item $2x+5y=6$
    \end{enumerate}
  \item
    \begin{enumerate}
      \item $\vx=s\colvecThree{2}{1}{2}+t\colvecThree{-3}{2}{1}$
      \item $\casesThree{x=2s-3t}{y=s+2t}{z=2s+t}$
    \end{enumerate}
  \item
    \begin{enumerate}
      \item $\vx=\colvecThree{6}{-4}{-3}+s\colvecThree{0}{1}{1}+t\colvecThree{-1}{1}{1}$
      \item $\casesThree{x=6-t}{y=-4+s+t}{z=-3+s+t}$
    \end{enumerate}
  \item $\vx=\colvecTwo{1}{-2}+t\colvecTwo{2}{2}$
  \item $\vx=\colvecThree{0}{1}{-1}+t\colvecThree{-2}{0}{4}$
  \item $\vx=\colvecThree{1}{1}{1}+s\colvecThree{3}{-1}{1}+t\colvecThree{-1}{0}{-2}$
  \item $\vx=\colvecThree{1}{1}{0}+s\colvecThree{0}{-1}{1}+t\colvecThree{-1}{0}{1}$
  \item
    \begin{enumerate}
      \item $\vx=\colvecTwo{0}{-1}+t\colvecTwo{1}{3}$, $\casesTwo{x=t}{y=-1+3t}$
      \item $\vx=\colvecTwo{0}{\frac{5}{2}}+t\colvecTwo{10}{-15}$,
        $\casesTwo{x=10t}{y=\frac{5}{2}-15t}$
    \end{enumerate}
  \item
    \begin{enumerate}
      \item $\mathbf{q}-\mathbf{p}$ is the directional vector. Thus, $\vx=\mathbf{p}+
        t(\mathbf{q}-\mathbf{p})$ is the vector form of the equation. Because $t$ varies
        from 0 to 1, the line segment is from $\mathbf{p}$ to $\mathbf{q}$.
      \item $\frac{1}{2}$. $\vx=\frac{1}{2}(\mathbf{p}+\mathbf{q})$.
      \item $(1, -1)$
      \item $(\frac{5}{2}, \frac{1}{2}, -\frac{1}{2})$
      \item $(\frac{4}{3}, -\frac{5}{3}), (\frac{2}{3}, -\frac{1}{3})$
    \end{enumerate}
  \item Direction vectors for the two lines are $\mathbf{d_1}=\colvecTwo{1}{m_1}$ and
    $\mathbf{d_2}=\colvecTwo{1}{m_2}$ The lines are perpendicular if and only if
    $\mathbf{d_1}$ and $\mathbf{d_2}$ are orthogonal. But $\mathbf{d_1}\cdot
    \mathbf{d_2}=0$ if and only if $1+m_1m_2=0$ or, equivalently, $m_1m_2=-1$.
  \item
    \begin{enumerate}
      \item Perpendicular.
      \item Parallel.
      \item Parallel.
      \item Perpendicular.
    \end{enumerate}
  \item
    \begin{enumerate}
      \item Perpendicular.
      \item Parallel.
      \item Perpendicular.
      \item Perpendicular.
    \end{enumerate}
  \item $\vx=\colvecTwo{2}{-1}+t\colvecTwo{2}{-3}$
  \item $\vx=\colvecTwo{2}{-1}+t\colvecTwo{-3}{2}$
  \item $\vx=\colvecThree{-1}{0}{3}+t\colvecThree{1}{-3}{2}$
  \item $\vx=\colvecThree{-1}{0}{3}+t\colvecThree{-1}{3}{-1}$
  \item $\colvecThree{6}{-1}{2}\cdot\vx=12$
  \item
    \begin{enumerate}
      \item $x=0, x=1, y=0, y=1, z=0, z=1$
      \item The plane is perpendicular to the $xy$-plane whose normal is $(0, 0, 1)$; thus, the
        plane's normal must have the form $(a, b, 0)$. The plane contains the diagonal
        $\vx=(1, 1, 1)$. $\mathbf{n}\cdot\vx=0\implies(a, b, 0)\cdot(1, 1, 1)=a+b+0=0 \implies
        b=-a$. Thus, $n$ can be $(1, -1, 0)$. The plane passes through the origin, thus its equation
        is $\vn\cdot\vx=0$, which is $x-y=0$.
      \item $x+y-z=0$.
    \end{enumerate}
  \item $2x+y+3z=10$
  \item $\frac{3\sqrt{2}}{2}$
  \item $\frac{5\sqrt{13}}{13}$
  \item $\frac{2\sqrt{3}}{3}$
  \item $\frac{1}{3}$
  \item $\left(\frac{1}{2}, \frac{1}{2}\right)$
  \item $\left(\frac{15}{13}, 1, \frac{10}{13}\right)$
  \item $\left(\frac{4}{3}, \frac{4}{3}, \frac{8}{3}\right)$
  \item $\left(\frac{10}{9}, \frac{7}{9}, \frac{11}{9}\right)$
  \item $\frac{18\sqrt{13}}{13}$
  \item $\frac{\sqrt{42}}{3}$
  \item $\frac{5}{3}$
  \item $\frac{2\sqrt{3}}{3}$
  \item Suppose a point $A(x_1, y_1)$ is on line $\ell$. This means $ax_1+by_1=c$. The normal vector
  of the line is $\vn=(a, b)$. Thus, the distance from the point B to the line is the length of the
  projection of $\overrightarrow{AB}$ on $\vn$.
  \item Same as in Exercise 39.
  \setcounter{enumi}{42}
  \item $\approx78.9^\circ$
  \item $\approx79.1^\circ$
  \item $\approx9.59^\circ$
  \setcounter{enumi}{46}
  \item $v-\operatorname{proj}_{\vn}(\vv)$
  \item \begin{enumerate}
    \item $\left(\frac{4}{3}, \frac{1}{3}, \frac{-5}{3}\right)$
  \end{enumerate}
\end{enumerate}

\subsection*{Exploration: The Cross Product}
\begin{enumerate}
  \item \begin{enumerate}
    \item $(3, 3, -3)$
    \item $(-3, -3, 3)$
  \end{enumerate}
 \setcounter{enumi}{3} 
 \item \begin{enumerate}
  \item $\colvecThree{3}{3}{-3}\cdot\colvecThree{x}{y}{z}=9$
  \item $\colvecThree{-5}{5}{5}\cdot\colvecThree{x}{y}{z}=0$
 \end{enumerate}
\end{enumerate}

\subsection{Applications}
\begin{enumerate}
  \item $(5, 12)$
  \item $(-15, -20)$
  \item A force with magnitude of $8\sqrt{3}$ N at an angle of $30^\circ$ to $\mathbf{f}_1$.
  \item Magnitude: $52-24\sqrt{2}$ N; angle: $\approx13.6^\circ$ to $\mathbf{f}_1$.
  \item Magnitude: 4; angle: $120^\circ$ to $\mathbf{f}_1$.
  \item Magnitude: $5\sqrt{2}$; angle: $45^\circ$ to $\mathbf{f}_1$.
  \item $(5, 5\sqrt{3})$
  \item 5
  \item $750\sqrt{2}$
  \item $50\sqrt{2}$
  \item 980 N
  \item $\frac{9.8\sqrt{2}}{2}$ N
  \item $\|\mathbf{f}_1\|\approx117.60$ N, $\|\mathbf{f}_2\|\approx88.20$ N
  \item $\|\mathbf{f}_1\|\approx143.48$ N, $\|\mathbf{f}_2\|\approx175.73$ N
\end{enumerate}

\subsection*{Chapter Review}
\begin{enumerate}
  \item \begin{enumerate}
    \item True
    \item False
    \item False
    \item False
    \item True
    \item True
    \item False
    \item True
    \item True
    \item False
  \end{enumerate}
  \item $(9, 12)$
  \item $(2, 7)$
  \item We have $\overrightarrow{BC} = \overrightarrow{OC}-\overrightarrow{OB}$.
  $\overrightarrow{OC}=-\overrightarrow{OA}$. Thus, $\overrightarrow{BC}=-\mathbf{a}-\mathbf{b}$.
  \item $120^\circ$
\end{enumerate}
\end{document}
